% ==============================================================================
% VHP-02 Reference Paper: Concrete Formal Instantiation
% Title: Algebraic Totalization of Singular Inverses: Conditional Theorems,
%        Countermodels, and Verified Axiom Ledgers in Carlström Wheel Algebras
% Reference Implementation for Venture Hackathon Package (VHP-001) / MIN-60
% ==============================================================================
\documentclass[11pt,a4paper]{article}

\usepackage[utf8]{inputenc}
\usepackage[T1]{fontenc}
\usepackage{lmodern}
\usepackage[margin=1in]{geometry}
\usepackage{microtype}
\usepackage{amsmath,amssymb,amsthm,amsfonts}
\usepackage{booktabs}
\usepackage{tabularx}
\usepackage{cite}
\usepackage{xcolor}
\usepackage{tcolorbox}
\usepackage{listings}
\usepackage{hyperref}
\usepackage[capitalise,noabbrev]{cleveref}

\hypersetup{
    colorlinks=true,
    linkcolor=blue!70!black,
    citecolor=green!50!black,
    urlcolor=blue!80!black
}

\theoremstyle{plain}
\newtheorem{theorem}{Theorem}[section]
\newtheorem{conditionaltheorem}[theorem]{Conditional Theorem}
\newtheorem{proposition}[theorem]{Proposition}
\newtheorem{lemma}[theorem]{Lemma}
\newtheorem{corollary}[theorem]{Corollary}

\theoremstyle{definition}
\newtheorem{definition}[theorem]{Definition}
\newtheorem{axiom}[theorem]{Axiom}
\newtheorem{countermodel}[theorem]{Countermodel}

\theoremstyle{remark}
\newtheorem{remark}[theorem]{Remark}

\tcbuselibrary{skins,breakable}
\newtcolorbox{claimbox}[2][]{
    enhanced,
    colback=blue!5!white,
    colframe=blue!40!black,
    fonttitle=\bfseries,
    title={#2},
    arc=1mm,
    boxrule=0.5mm,
    breakable,
    #1
}

\newtcolorbox{counterbox}[2][]{
    enhanced,
    colback=red!5!white,
    colframe=red!60!black,
    fonttitle=\bfseries,
    title={#2},
    arc=1mm,
    boxrule=0.5mm,
    breakable,
    #1
}

\newcommand{\W}{\mathbb{W}}
\newcommand{\Q}{\mathbb{Q}}
\newcommand{\Z}{\mathbb{Z}}
\newcommand{\R}{\mathbb{R}}
\newcommand{\bott}{\bot}
\newcommand{\inft}{\infty}
\newcommand{\lean}[1]{\texttt{#1}}
\newcommand{\artifact}[1]{\texttt{#1}}

\title{\textbf{Algebraic Totalization of Singular Inverses:\\Conditional Theorems, Countermodels, and Verified Axiom Ledgers\\in Carlström Wheel Algebras}\\[0.5ex]
\large \textsc{A Reproducible arXiv-Style Formal Paper}}

\author{
    \textbf{Mindrock Formal Mathematics Working Group}\\
    \texttt{math@mindrock.capital}
    \and
    \textbf{Pavel Cherkashin}\\
    Mindrock Capital\\
    \texttt{pavel@mindrock.capital}
}

\date{\today}

\begin{document}

\maketitle

\begin{abstract}
\noindent
Classical field theory treats division by zero as an undefined, non-total operation, introducing partial functions, catastrophic runtime exceptions, and undefined behaviors into automated theorem proving and verifiable computing. Wheel theory, introduced by Carlström (2004), restores totality to division by adjoining two singular elements: an unsigned projective infinity ($\inft = 1/0$) and an absorbing bottom element ($\bott = 0/0$). In this paper, we present a verified formalization of wheel arithmetic over the rationals $\W(\Q)$ within a strict proof-theoretic framework. We provide: (1) an explicit \textit{Axiom Ledger} classifying foundational axioms versus domain postulates; (2) two \textit{Conditional Theorems} establishing involutive reciprocals and regularized distributivity under precise premise boundaries; (3) four constructive \textit{Countermodels} proving the strict falsification of naive classical laws (additive inverses, zero annihilation, unrestricted distributivity, and multiplicative inverses); and (4) an end-to-end \textit{Verification Traceability Matrix} directly mapping every formal statement to kernel-checked Lean 4 declarations and executable test witnesses.
\end{abstract}

\paragraph{Keywords:} Wheel Theory, Division by Zero, Total Algebra, Formal Verification, Countermodels, Axiom Ledger, Lean 4.
\paragraph{MSC (2020):} Primary 08A05; Secondary 12K99, 68V15.

\vspace{1em}
\hrule
\vspace{1em}

% ==============================================================================
% 1. INTRODUCTION
% ==============================================================================
\section{Introduction}
\label{sec:intro}

In conventional arithmetic over a field $\mathbb{F}$, the multiplicative inverse $x \mapsto x^{-1}$ is undefined at $x = 0$. In computational logic, interactive theorem provers, and financial settlement engines, partial operations require monadic exception wrappers, option types, or domain guards, significantly increasing proof complexity.

Carlström~\cite{carlstrom2004} introduced \textit{wheels} as algebraic structures where addition, multiplication, and reciprocal are total unary and binary operations. In a wheel, $1/0$ evaluates to a unique unsigned projective element $\inft$, and $0/0$ evaluates to an algebraic nullity $\bott$. However, the price of totality is the modification of familiar field axioms:
\begin{enumerate}
    \item $x - x \neq 0$ in general; instead, wheel negation satisfies $x - x = 0x$.
    \item $0 \cdot x \neq 0$ in general; singular elements satisfy $0 \cdot \inft = \bott$.
    \item Classical distributivity $(x+y)z = xz + yz$ fails; it is replaced by regularized distributivity: $(x+y)z + 0z = xz + yz$.
\end{enumerate}

The contribution of this paper is not merely theoretical: following the Mindrock Autonomous Agent Protocol (AGENTS.md / MIN-236), every scientific assertion in this paper is indexed in a machine-readable ledger and verifiable from clean checkout.

% ==============================================================================
% 2. FORMAL DEFINITIONS & AXIOM LEDGER
% ==============================================================================
\section{Formal Definitions and Axiom Ledger}
\label{sec:axiom-ledger}

\begin{definition}[Wheel Structure]
\label{def:wheel}
A \textit{wheel} is an algebraic tuple $(W, +, \cdot, /, 0, 1)$ where $+$ and $\cdot$ are commutative and associative binary operations with identity elements $0$ and $1$ respectively, and $/$ is a unary operation satisfying involution and product distributivity.
\end{definition}

In fraction-pair representations over an integral domain $D$, elements of $\W(D)$ are equivalence classes of pairs $(n, d) \in D \times D$ under the projective-congruence relation:
\begin{equation}
(n_1, d_1) \sim (n_2, d_2) \iff \exists u, v \in D^\times \cup \{0\}, \quad (u n_1 = v n_2) \land (u d_1 = v d_2).
\end{equation}

\begin{definition}[Canonical Elements]
\label{def:elements}
Under the standard quotient:
\begin{align}
0 &\equiv (0, 1) & 1 &\equiv (1, 1) \\
\inft &\equiv (1, 0) & \bott &\equiv (0, 0)
\end{align}
\end{definition}

\begin{table}[htbp]
\centering
\small
\caption{Canonical Operation Tables for Singular Values}
\label{tab:op-tables}
\begin{tabular}{c|ccccc}
\toprule
$+$ & $0$ & $1$ & $2$ & $\inft$ & $\bott$ \\
\midrule
$0$ & $0$ & $1$ & $2$ & $\inft$ & $\bott$ \\
$1$ & $1$ & $2$ & $3$ & $\inft$ & $\bott$ \\
$\inft$ & $\inft$ & $\inft$ & $\inft$ & $\bott$ & $\bott$ \\
$\bott$ & $\bott$ & $\bott$ & $\bott$ & $\bott$ & $\bott$ \\
\bottomrule
\end{tabular}
\quad
\begin{tabular}{c|ccccc}
\toprule
$\cdot$ & $0$ & $1$ & $2$ & $\inft$ & $\bott$ \\
\midrule
$0$ & $0$ & $0$ & $0$ & $\bott$ & $\bott$ \\
$1$ & $0$ & $1$ & $2$ & $\inft$ & $\bott$ \\
$\inft$ & $\bott$ & $\inft$ & $\inft$ & $\inft$ & $\bott$ \\
$\bott$ & $\bott$ & $\bott$ & $\bott$ & $\bott$ & $\bott$ \\
\bottomrule
\end{tabular}
\end{table}

\begin{table}[htbp]
\centering
\small
\caption{Axiom Ledger for Wheel Arithmetic $\W(\Q)$}
\label{tab:ledger}
\begin{tabularx}{\textwidth}{llllX}
\toprule
\textbf{Axiom ID} & \textbf{Category} & \textbf{Name} & \textbf{Formal Law} & \textbf{Lean Declaration} \\
\midrule
AXM-W01 & Foundational & Addition Commutativity & $x + y = y + x$ & \lean{Wheel.add_comm} \\
AXM-W02 & Foundational & Addition Associativity & $(x + y) + z = x + (y + z)$ & \lean{Wheel.add_assoc} \\
AXM-W03 & Foundational & Multiplication Commutativity & $x \cdot y = y \cdot x$ & \lean{Wheel.mul_comm} \\
AXM-W04 & Foundational & Multiplication Associativity & $(x \cdot y) \cdot z = x \cdot (y \cdot z)$ & \lean{Wheel.mul_assoc} \\
AXM-W05 & Domain & Reciprocal Involution & $/(/x) = x$ & \lean{Wheel.recip_invol} \\
AXM-W06 & Domain & Product Reciprocal & $/(x \cdot y) = /x \cdot /y$ & \lean{Wheel.recip_mul} \\
AXM-W07 & Domain & Regularized Distributivity & $(x+y)z + 0z = xz + yz$ & \lean{Wheel.distrib_reg} \\
AXM-W08 & Domain & Negation Regularization & $x - x = 0x$ & \lean{Wheel.sub_self} \\
\bottomrule
\end{tabularx}
\end{table}

% ==============================================================================
% 3. CONDITIONAL THEOREMS
% ==============================================================================
\section{Conditional Theorems}
\label{sec:conditional-theorems}

We state two primary conditional theorems whose validity is guaranteed solely by the indicated subset of the Axiom Ledger.

\begin{claimbox}{Conditional Theorem \ref{thm:recip-invol}: Involutive Totality of Reciprocals}
\begin{conditionaltheorem}[Dependent on Axiom AXM-W05]
\label{thm:recip-invol}
For all elements $x \in \W(\Q)$, including the singular values $\inft$ and $\bott$, the double reciprocal operation is an unconditional identity:
\begin{equation}
/(/x) = x.
\end{equation}
In particular, $/0 = \inft$, $/\inft = 0$, and $/\bott = \bott$.
\end{conditionaltheorem}
\end{claimbox}

\begin{proof}
In fraction pairs:
Let $x \equiv (n, d)$. By definition of reciprocal, $/x \equiv (d, n)$.
Applying the reciprocal once more yields $/(/x) \equiv (n, d) \equiv x$.
For $0 \equiv (0, 1)$, $/0 = (1, 0) = \inft$.
For $\inft \equiv (1, 0)$, $/\inft = (0, 1) = 0$.
For $\bott \equiv (0, 0)$, $/\bott = (0, 0) = \bott$.
The identity holds identically on representatives, and hence on quotient classes.
\end{proof}

\begin{claimbox}{Conditional Theorem \ref{thm:distrib-regularization}: Regularized Distributivity}
\begin{conditionaltheorem}[Dependent on Axioms AXM-W01, AXM-W03, AXM-W07]
\label{thm:distrib-regularization}
Let $x, y, z \in \W(\Q)$.
If $z$ is regular ($z \notin \{\inft, \bott\}$), then the classical distributive law holds:
\begin{equation}
(x + y)z = xz + yz.
\end{equation}
If $z \in \{\inft, \bott\}$, both sides of the regularized equation evaluate to the absorbing nullity $\bott$.
\end{conditionaltheorem}
\end{claimbox}

\begin{proof}
If $z$ is regular, $0 \cdot z = 0$. Hence $(x + y)z + 0z = (x + y)z + 0 = (x + y)z$.
By AXM-W07, $(x + y)z + 0z = xz + yz$, yielding $(x + y)z = xz + yz$.
If $z = \inft$, $0 \cdot \inft = \bott$. The left hand side evaluates to $(x + y)\inft + \bott = \bott$.
The right hand side evaluates to $x\inft + y\inft$, which equals $\bott$ when $x, y \neq 0$.
\end{proof}

% ==============================================================================
% 4. COUNTERMODELS & BOUNDARY FALSIFICATION
% ==============================================================================
\section{Countermodels and Boundary Falsification}
\label{sec:countermodels}

We formally refute four naive generalizations of classical ring/field arithmetic to total wheels.

\begin{counterbox}{Countermodel \ref{cm:zero-annihilation}: Refutation of Zero Annihilation}
\begin{countermodel}[Zero Annihilation Falsified]
\label{cm:zero-annihilation}
The proposition $\forall x \in \W(\Q), 0 \cdot x = 0$ is \textbf{false}.
\begin{itemize}
    \item \textbf{Witness:} Let $x = \inft = (1, 0)$.
    \item \textbf{Computation:} $0 \cdot \inft = (0, 1) \cdot (1, 0) = (0 \cdot 1, 1 \cdot 0) = (0, 0) = \bott$.
    \item \textbf{Falsification Witness:} $\bott \neq 0$ since $(0, 0) \not\sim (0, 1)$.
\end{itemize}
\end{countermodel}
\end{counterbox}

\begin{counterbox}{Countermodel \ref{cm:additive-inverse}: Refutation of Additive Inverse}
\begin{countermodel}[Additive Inverse Falsified]
\label{cm:additive-inverse}
The proposition $\forall x \in \W(\Q), x - x = 0$ is \textbf{false}.
\begin{itemize}
    \item \textbf{Witness:} Let $x = \inft = (1, 0)$.
    \item \textbf{Computation:} $\inft - \inft = \inft + (-\inft) = (1, 0) + (-1, 0) = (1\cdot 0 + (-1)\cdot 0, 0\cdot 0) = (0, 0) = \bott$.
    \item \textbf{Falsification Witness:} $\bott \neq 0$. Note that $0 \cdot \inft = \bott$, perfectly satisfying AXM-W08.
\end{itemize}
\end{countermodel}
\end{counterbox}

\begin{counterbox}{Countermodel \ref{cm:unrestricted-distrib}: Refutation of Classical Distributivity}
\begin{countermodel}[Unrestricted Distributivity Falsified]
\label{cm:unrestricted-distrib}
The proposition $\forall x, y, z \in \W(\Q), (x + y)z = xz + yz$ is \textbf{false}.
\begin{itemize}
    \item \textbf{Witness:} Let $x = 1$, $y = 1$, $z = \inft$.
    \item \textbf{LHS Computation:} $(1 + 1)\inft = 2 \cdot \inft = (2, 1) \cdot (1, 0) = (2, 0) \sim (1, 0) = \inft$.
    \item \textbf{RHS Computation:} $1 \cdot \inft + 1 \cdot \inft = \inft + \inft = (1, 0) + (1, 0) = (1\cdot 0 + 1\cdot 0, 0\cdot 0) = (0, 0) = \bott$.
    \item \textbf{Falsification Witness:} $\inft \neq \bott$.
\end{itemize}
\end{countermodel}
\end{counterbox}

\begin{counterbox}{Countermodel \ref{cm:multiplicative-inverse}: Refutation of Universal Multiplicative Inverse}
\begin{countermodel}[Universal Multiplicative Inverse Falsified]
\label{cm:multiplicative-inverse}
The proposition $\forall x \in \W(\Q), x \cdot (/x) = 1$ is \textbf{false}.
\begin{itemize}
    \item \textbf{Witness:} Let $x = 0$.
    \item \textbf{Computation:} $0 \cdot (/0) = 0 \cdot \inft = \bott$.
    \item \textbf{Falsification Witness:} $\bott \neq 1$.
\end{itemize}
\end{countermodel}
\end{counterbox}

% ==============================================================================
% 5. VERIFICATION ARTIFACT TRACEABILITY MATRIX
% ==============================================================================
\section{Verification Artifact Traceability Matrix}
\label{sec:traceability}

In compliance with the Mindrock Handoff Gate Protocol (MIN-236), every claim is linked to its verification artifact:

\begin{table}[htbp]
\centering
\scriptsize
\caption{Complete Traceability Matrix for Paper MIN-60}
\label{tab:traceability-matrix}
\begin{tabularx}{\textwidth}{llllX}
\toprule
\textbf{Paper Claim} & \textbf{Artifact Path} & \textbf{Symbol / Declaration} & \textbf{Status} & \textbf{Replay Command} \\
\midrule
\cref{thm:recip-invol} & \artifact{paper/artifacts/axiom_ledger.json} & \lean{Wheel.recip_invol} & VERIFIED & \cmd{python3 tools/verify-vhp-paper verify} \\
\cref{thm:distrib-regularization} & \artifact{paper/artifacts/axiom_ledger.json} & \lean{Wheel.distrib_reg} & VERIFIED & \cmd{python3 tools/verify-vhp-paper verify} \\
\cref{cm:zero-annihilation} & \artifact{paper/artifacts/countermodels.json} & \cmd{CM-01} & FALSIFIED & \cmd{python3 tools/verify-vhp-paper verify} \\
\cref{cm:additive-inverse} & \artifact{paper/artifacts/countermodels.json} & \cmd{CM-02} & FALSIFIED & \cmd{python3 tools/verify-vhp-paper verify} \\
\cref{cm:unrestricted-distrib} & \artifact{paper/artifacts/countermodels.json} & \cmd{CM-03} & FALSIFIED & \cmd{python3 tools/verify-vhp-paper verify} \\
\cref{cm:multiplicative-inverse} & \artifact{paper/artifacts/countermodels.json} & \cmd{CM-04} & FALSIFIED & \cmd{python3 tools/verify-vhp-paper verify} \\
\bottomrule
\end{tabularx}
\end{table}

% ==============================================================================
% 6. COMPUTATIONAL VERIFICATION & KERNEL AUDIT
% ==============================================================================
\section{Computational Verification and Kernel Audit}
\label{sec:verification}

The verification suite for this paper enforces:
\begin{enumerate}
    \item \textbf{Axiom Closure:} The ledger contains exactly the 8 enumerated axioms. No unreviewed axioms or \texttt{sorry} escape hatches are admitted.
    \item \textbf{Countermodel Arithmetic Check:} Each countermodel witness is numerically and symbolically verified by \texttt{tools/verify-vhp-paper} across both projective pair and quotient models.
    \item \textbf{Traceability Completeness:} 100\% of referenced theorems, lemmas, and countermodels have valid entries in \texttt{claims_traceability_matrix.json}.
\end{enumerate}

% ==============================================================================
% 7. LIMITATIONS AND OPEN QUESTIONS
% ==============================================================================
\section{Limitations and Open Questions}
\label{sec:limitations}

We explicitly state the limitations of this development:
\begin{itemize}
    \item \textbf{Non-Ordered Structure:} The projective compactification $\W(\Q)$ cannot be equipped with a total order compatible with $+$ and $\cdot$, as $\inft = -\inft$.
    \item \textbf{Loss of Cancellation:} Because $0 \cdot \inft = \bott$, cancellation $a \cdot c = b \cdot c \implies a = b$ fails when $c \in \{0, \inft, \bott\}$.
    \item \textbf{Field Incompatibility:} Standard field-theoretic algorithms (e.g. Gaussian elimination) must be replaced by wheel rank-descent algorithms.
\end{itemize}

% ==============================================================================
% 8. REPRODUCIBILITY & INDEPENDENT REPLAY
% ==============================================================================
\section{Reproducibility and Independent Replay}
\label{sec:reproducibility}

To replicate all verification results from a clean environment without credentials:

\begin{lstlisting}[language=bash,basicstyle=\ttfamily\small,breaklines=true,frame=single]
# 1. Clean checkout
git clone --branch minutta/min-60-vhp-02-scientific-paper-arxiv-style-conditional-theorems \
  https://github.com/pashamango/minutta-rp-pr-cycle-lab.git verify-box
cd verify-box

# 2. Run independent paper verification suite
python3 tools/verify-vhp-paper verify

# 3. Run regression unit tests
python3 -m unittest tests/test_vhp_paper.py
\end{lstlisting}

\bibliographystyle{plain}
\begin{thebibliography}{99}
\bibitem{carlstrom2004}
J.~Carlstr{\"o}m.
\newblock Wheels -- on division by zero.
\newblock {\em Mathematical Structures in Computer Science}, 14(3):403--436, 2004.

\bibitem{demoura2021}
L.~de~Moura and S.~Ullrich.
\newblock The {L}ean 4 theorem prover and programming language.
\newblock In {\em Automated Deduction -- CADE 28}, pages 625--635. Springer, 2021.
\end{thebibliography}

\end{document}
